\section{Independent degree-summation arguments}
\label{l:degree-supplement}

The quantitative toolkit in Section~\ref{l:degrees} proves every
inverse-power limit needed by the transfers. The arguments below retain
distinct ways to obtain those limits: some use only the composition
bound on partitions, and some avoid the hook formula. We use the same
notation $L(\lambda)$, $k(\lambda)$, $D_\lambda$, $C_u$ and $Z_u(n)$
as in that section. The common tableau and hook estimates are
Lemmas~\ref{l:tableau-constructions} and~\ref{l:degree-hook-ratio}.

\subsection{A hook proof at the twentieth power}

The following proof needs only the composition bound on the number of
partitions.  Its explicit exponent is useful when an argument counts
small subgroup types without the sharper partition estimate.

\begin{lemma}[Inverse-twentieth powers and type counts]
\label{l:degree-reciprocal-twenty}
One has $Z_{20}(n)\to0$.  There is an absolute constant $K_0$ such
that, for all $n\ge1$ and $R\ge1$,
\begin{equation}\label{l:degree-twenty-type-count}
 \#\{\lambda\vdash n:D_\lambda\le R\}\le K_0R^{20}.
\end{equation}
\end{lemma}
\begin{proof}
Orient the first row to have length $L=L(\lambda)$ and put $k=n-L$.
If $L<n/8$, then $D_\lambda\ge(4/e)^n$, so the contribution is at
most $[2(e/4)^{20}]^n=o(1)$.
For $L\ge n/8$ and $k\ge1$, \eqref{l:hook-rearrangement} gives
\[
 D_\lambda\ge\binom nk
                  \exp\left(-\frac{k(1+\log L)}L\right).
\]
There are at most $2\cdot2^k$ shapes at each $k$.
For $1\le k\le\sqrt n$, the exponent is uniformly $o(1)$ and
$\binom nk\ge(n/k)^k\ge n^{k/2}$.  The contribution is at most
\[
 2e^{o(1)}\sum_{k=1}^{\lfloor\sqrt n\rfloor}(2n^{-10})^k=o(1).
\]
For $\sqrt n<k\le7n/8$, use
$\binom nk\ge(8/7)^k$ and
$k(1+\log L)/L\le7(1+\log n)$.  This contribution is at most
\[
 2e^{140}n^{140}
       \sum_{k>\sqrt n}[2(7/8)^{20}]^k=o(1),
\]
since $2(7/8)^{20}<1$.  The full inverse-twentieth sum is uniformly
bounded after adjoining the row and column and the finitely many small
sizes.  Each type of dimension at most $R$ contributes at least
$R^{-20}$ to that sum, proving \eqref{l:degree-twenty-type-count}.
\end{proof}

\subsection{Direct reciprocal and inverse-square proofs}
\label{l:degree-direct-routes}

These proofs share a summation step, not a dimension estimate.
At deficit $k\ge1$ there are at most $2p(k)$ shapes. A lower bound
that diverges at each fixed $k$ and has a summable inverse-power
majorant therefore gives a vanishing near-row and near-column sum.
In the remaining ranges, $D_\lambda\ge e^{cn}$, or even
$e^{cn^{0.6}}$, beats $p(n)\le e^{3\sqrt n}$.
Adding the row and column and finitely many small sizes then gives
the corresponding constant $C_u$. Each argument below specifies
its dimension bounds and the partition count it uses.

\begin{lemma}[Direct inverse-square summability]
\label{l:degree-inverse-square}\label{l:degree-reciprocal-two}
One has $C_2<\infty$ and $Z_2(n)\to0$. This conclusion admits
each of the tableau, hook-ratio, and second-row proofs below.
\end{lemma}
\begin{proof}[Initial-row interleavings and diagonals]
Orient a longest line as the first row $L=n-k$.
For $1\le k\le n/3$, the bound
$D_\lambda\ge\binom{n-k}{k}\ge2^k$ and the composition count
$2p(k)\le2\cdot2^k$ give the majorant $2\cdot2^{-k}$,
with fixed-$k$ divergence. If $k>n/3$ and $L\ge n/8$, retain
$\lfloor L/2\rfloor$ lower boxes to obtain
\[
 D_\lambda\ge\binom L{\lfloor L/2\rfloor}
 \ge\frac{2^L}{L+1}\ge\frac{2^{n/8}}{n+1}.
\]
If $L<n/8$, diagonal filling gives $D_\lambda\ge2^{3n/4}$.
The shared summation step completes this hook-free proof.

The alternative diagonal cutoff $L<n/4$ gives $D_\lambda\ge2^{n/2}$.
For $L\ge n/4$, retain $b=\min(k,\lfloor L/2\rfloor)$:
$k\le L/2$ gives the same $2^k$ majorant, while $k>L/2$
gives the exponential central-binomial bound.
\end{proof}

\begin{proof}[Initial-row interleavings with factorial or hook tails]
At $k\le n/4$, use $\binom{n-k}k\ge3^k$, with inverse-square
majorant $2(2/9)^k$; at $k\le n/3$, use $2^k$ and
$2\cdot2^{-k}$. Either binomial diverges for fixed positive $k$.
The following rows give alternative completions of these near-row
bounds. In each middle range retain the first row and $b$ lower boxes:
\[
\begin{array}{c|c|c}
 \text{middle range}&b&\text{lower bound for }D_\lambda\\ \hline
 k>n/4,\ L\ge n/10&\lfloor L/4\rfloor&4^b\\
 k>n/4,\ L\ge n/8&\lfloor L/4\rfloor&4^b\\
 k>n/4,\ L\ge n/8&\lfloor n/32\rfloor&4^b\\
 k>n/4,\ L\ge n/8&\lfloor n/16\rfloor&2^b\\
 k>n/3,\ L\ge n/10&\lfloor L/3\rfloor&3^b\\
 k>n/4,\ L\ge n^{0.6}&\lfloor L/4\rfloor&4^b
\end{array}
\]
Every row follows from $D_\lambda\ge\binom Lb\ge(L/b)^b$.
There are enough lower boxes: $k>n/4$ implies $k>L/3$,
and $k>n/3$ implies $k>L/2$; the stated cutoffs also ensure
$b\le L$. The fixed-cutoff rows are exponential in $n$.
Below $L=n/10$ or $n/8$, the factorial form of diagonal filling
gives, respectively, $(5/e)^n$ or $(4/e)^n$.
These are hook-free completions; using the hook formula gives
the same two bounds independently. For the moving cutoff,
the middle bound is $e^{cn^{0.6}}$ and hooks below the cutoff give
$n!/(2n^{0.6})^n$. Thus every row completes the inverse-square sum.
\end{proof}

\begin{proof}[First-row hook ratios]
For $k\le n/4$, \eqref{l:hook-short-tail} gives
$D_\lambda\ge e^{-1}\binom nk$, so the inverse-square sum is at most
\[
 2e^2\sum_{1\le k\le n/4}2^k(k/n)^{2k}=o(1).
\]
The summands inside the sum are at most $8^{-k}$ and vanish at
fixed $k$. In the remainder $L<3n/4$, use $(5/e)^n$ for
$L\le n/10$; otherwise retain $b=\lfloor L/4\rfloor$ lower boxes.
The hook ratio gives $D_\lambda\ge e^{-1}\binom{L+b}b\ge e^{-1}4^b$,
where $b\ge n/50$ for large $n$.

Two alternative hook cutoffs retain different quantitative losses.
For $L\ge n/2$, opposite rearrangement gives
\[
 D_\lambda\ge\binom nk e^{-k(1+\log L)/L}
 \ge2^k e^{-k(1+\log L)/L}\ge(\sqrt2)^k
\]
for large $n$. Here use $2p(k)=e^{O(\sqrt k)}$ for domination;
the binomial gives fixed-$k$ divergence. For $n/10\le L<n/2$,
$\binom nL$ is exponential in $n$ and the logarithmic correction
is $O(\log n)$; below $n/10$ use $(5/e)^n$.
Alternatively, split at $L=3n/4$ and $n/8$: the first range has
short-tail loss $k/(L-k+1)\le1$, the middle uses the same
$b=\lfloor L/4\rfloor$ hook subdiagram, and the last uses $(4/e)^n$.
Each choice proves the limit by the shared summation argument.
\end{proof}

\begin{proof}[Second-row width and two-row rectangles]
Write $c=\lambda_2$ after orientation. At $1\le k\le n/4$,
\[
 D_\lambda\ge\binom{n-c}k\ge\binom{n-k}k.
\]
Its ratio base is at least $\sqrt n/2$ for $k\le\sqrt n$
and at least three for $\sqrt n<k\le n/4$, so the composition
count gives vanishing geometric sums. In the remainder, hooks
handle $L<n/10$. For $L\ge n/10$ and $c\le L/2$,
\[
 D_\lambda\ge\binom{k+L-c}k
 \ge2^{\min(k,L-c)}\ge2^{n/20}.
\]
The middle inequality follows by retaining $\min(k,L-c)$ factors
in the binomial product. If $c>L/2$, the diagram contains a
two-row rectangle of width $\lfloor n/20\rfloor$, whose Catalan
number is exponential in $n$. The subexponential total count finishes.
The alternative near-row cutoff $k\le n/3$ gives majorant
$2\cdot2^{-k}$; in its remainder $n/10\le L<2n/3$,
the same minimum is at least $n/20$, and the rectangle may have width $c$.
\end{proof}

\begin{proof}[Direct reciprocal bounds by interleaving and ballot words]
Interleaving gives $D_\lambda\ge\binom{n-k}k\ge3^k$
on $1\le k\le n/4$. Use the sharper deficit count
$2p(k)\le2e^{3\sqrt k}$ to sum $2p(k)3^{-k}$.
For $k>n/4$, hooks give $2^n$ if $L<n/(4e)$; otherwise
$b=\min(k,\lfloor L/2\rfloor)$ is a positive linear fraction
of $n$ and $\binom Lb\ge2^b$. Equivalently, split at $L=n/8$,
using $b=\lfloor n/16\rfloor$ above and $(4/e)^n$ below.
Both give $C_1<\infty$ and $Z_1(n)\to0$ without a near-row hook ratio.

Ballot words permit the coarser composition count near a long row.
For $L\ge n/2$, the binomial difference in \eqref{l:tableau-ballot}
diverges at fixed $k$, while its Catalan bound is at least
$4^k/((k+1)(2k+1))$. The reciprocal contribution is therefore
dominated by a constant times $(k+1)(2k+1)2^{-k}$.
For $n/10\le L<n/2$, retain $L$ lower boxes and use their Catalan
bound; below $n/10$, hooks give $(5/e)^n$.

A quarter-tail ballot completion instead splits at $L=n/8$ and
$3n/4$. The first range has hook bound $(4/e)^n$; in the middle,
$b=\lfloor L/4\rfloor$ lower boxes give
\[
 D_\lambda\ge\binom{L+b}b\frac{L-b+1}{L+1},
\]
exponential in $n$. For $L>3n/4$, the ballot fraction is at least
$1/2$, so $D_\lambda\ge\frac12\binom nk\ge\frac12 4^k$.
The deficit count and fixed-$k$ divergence finish this version as well.
\end{proof}

\begin{proof}[A reciprocal hook proof]
\phantomsection
\label{l:degree-small-index-reciprocal}
The hook estimates in Lemma~\ref{l:small-index-character} give
$D_\lambda\ge e^{n/64}$ for $n\ge64$ and $L\le3n/4$,
which beats the subexponential total count. For $L>3n/4$ and $k\ge1$,
\[
 D_\lambda\ge e^{-1}(n/k)^k\ge e^{-1}4^k.
\]
The composition count gives the reciprocal majorant $2e\,2^{-k}$,
and the first bound diverges at fixed $k$. Thus $Z_1(n)\to0$ and
$C_1<\infty$ follow independently from the hook comparisons used
to control the subgroup index.
\end{proof}

\subsection{Coarse partition counts and square-root dimension bounds}
\label{l:degree-square-root-section}

\begin{lemma}[Elementary sixteenth and thirteenth powers]
\label{l:degree-elementary-sixteen}
The conclusions $C_u<\infty$ and $Z_u(n)\to0$ for $u=13,16$
follow from tableau counting and $p(n)\le2^n$ alone.
\end{lemma}
\begin{proof}
If $1\le k\le n/3$, the initial-row construction gives
$D_\lambda\ge\binom{n-k}k\ge2^k$, with at most $2\cdot2^k$
shapes.  The contribution is bounded by $2\cdot2^{-(u-1)k}$,
and each fixed-$k$ contribution tends to zero.
For $k>n/3$ and $L\ge n/8$, retain
$\lfloor L/2\rfloor$ lower boxes to get
$D_\lambda\ge2^{n/8}/(n+1)$.
The remaining contribution in this range is at most
\[
 (n+1)^u2^{(1-u/8)n}=o(1).
\]
For $L<n/8$, diagonal filling gives $D_\lambda\ge2^{3n/4}$,
and the contribution is at most $2^{(1-3u/4)n}=o(1)$.
At $u=16$ these two bounds are exactly
$(n+1)^{16}2^{-n}$ and $2^{-11n}$.
At $u=13$ they vanish because $13/8>1$ and $39/4>1$.
The finitely many small-size sums give the uniform constants.
\end{proof}

\begin{lemma}[A tail-adapted diagonal bound]\label{l:degree-diagonal-four}
For all sufficiently large $n$ and every shape of deficit
$k=k(\lambda)\ge1$,
\[
 D_\lambda\ge2^{k/2},\qquad D_\lambda\ge L(\lambda)\ge\sqrt n.
\]
These inequalities and $p(k)\le2^k$ prove $C_4<\infty$ and
$Z_4(n)\to0$.
\end{lemma}
\begin{proof}
Let $a$ be the largest value of $i+j-1$ over cells $(i,j)$.
Diagonal filling gives $D_\lambda\ge2^{n-a}$.
Every cell satisfies $ij\le n$, so its smaller index is at most
$\sqrt n$ and $a\le L+\sqrt n$.
If $k<n/2$, orient the first row as $L$.  A lower cell satisfies
$(i-1)j\le k$, whence $i+j-1\le k+1\le L$; thus $a=L$.
If $k\ge n/2$, then $n-a\ge k-\sqrt n\ge k/2$ for large $n$.
This proves the first bound in both cases.
A non-row, non-column shape contains $(L,1)$ after orientation;
that hook has $L$ tableaux, proving the second bound.
For each $k$, at most $2p(k)\le2\cdot2^k$ shapes occur, so their
inverse fourth powers are bounded by $2\cdot2^{-k}$.
At fixed positive $k$ they tend to zero by $D_\lambda\ge\sqrt n$.
Dominated convergence and the finite-small-size argument finish.
\end{proof}

\begin{lemma}[Inverse-tenth powers from a square-root first line]
\label{l:degree-inverse-ten}
One has $C_{10}<\infty$ and $Z_{10}(n)\to0$, using only the
first-row interleaving construction and the bound $p(n)\le e^{3\sqrt n}$.
\end{lemma}
\begin{proof}
Orient the longest line as the first row $L\ge\sqrt n$.
For $1\le k<L/2$, use $D_\lambda\ge\binom Lk\ge2^k$;
the contribution at $k$ is at most $2\cdot2^{-9k}$ and tends to
zero for fixed $k$.
For $k\ge L/2$, retain $\lfloor L/2\rfloor$ lower boxes.
The bound $D_\lambda\ge2^L/(L+1)$ gives total contribution at most
\[
 e^{3\sqrt n}(n+1)^{10}2^{-10\sqrt n}=o(1).
\]
This proves the limit and the uniform bound.
\end{proof}

\paragraph{Central-binomial powers eight and twenty.}
Divide at $k\le n/3$. Initial-row interleaving gives
$D_\lambda\ge\binom{n-k}k\ge2^k$, with fixed-$k$ divergence
and near-row majorants $2\cdot2^{-7k}$ and $2\cdot2^{-19k}$.
For $k>n/3$, one has $k>L/2$ and $L\ge\sqrt n$.
Retaining $\lfloor L/2\rfloor$ lower boxes gives
$D_\lambda\ge2^L/(L+1)$, so the remaining sums are at most
\[
 e^{3\sqrt n}(n+1)^u2^{-u\sqrt n}=o(1)
 \qquad(u=8,20),
\]
because $8\log2>3$. This proves both limits directly from
central-binomial interleavings, without ballot words.

\paragraph{Weaker central-binomial estimates.}
For powers $16$ and $32$, the same near-row cutoff and
$\binom L{\lfloor L/2\rfloor}\ge2^{\lfloor L/2\rfloor}$ give
\begin{equation}\label{l:degree-sqrt-sixteen}
 e^{3\sqrt n}2^{-u\lfloor\sqrt n/2\rfloor}=o(1)
       \quad(u=16,32).
\end{equation}
Unlike Lemma~\ref{l:degree-elementary-sixteen}, these proofs use
only a square-root lower bound on $L$ and no diagonal filling.
For power $12$, divide at $k\le\lfloor L/2\rfloor$:
the near-row bound $\binom Lk\ge2^k$ gives majorant
$2\cdot2^{-11k}$ and fixed-$k$ divergence. In the remainder,
$\binom L{\lfloor L/2\rfloor}\ge2^{L/2}$ for large $L$ gives
$e^{3\sqrt n-6(\log2)\sqrt n}=o(1)$.

\paragraph{The ballot inverse-eighth route.}
For $1\le k\le L$, the ballot fraction is at least $1/(k+1)$.
Thus $D_\lambda\ge2^k/(k+1)$ gives summable majorant
$2^k(k+1)^8 2^{-8k}$, using $2p(k)\le2^k$; the full
binomial difference diverges at fixed $k$. For $k>L$, retain
$L$ lower boxes to get
$D_\lambda\ge\binom{2L}L/(L+1)\ge2^L/(L+1)$.
The remaining sum is at most
$e^{3\sqrt n}(n+1)^8 2^{-8\sqrt n}=o(1)$.
For every route, adjoining the two linear types and finitely many
small sizes gives the complete sum's uniform boundedness.

\begin{lemma}[A square-root lower bound in the deficit]
\label{l:degree-sqrt-deficit}
For every $n\ge1$ and $\lambda\vdash n$,
\begin{equation}\label{l:degree-sqrt-deficit-equation}
 \log D_\lambda\ge\frac{\log2}{2}\sqrt{k(\lambda)}.
\end{equation}
In particular $C_{32}<\infty$ follows from this bound alone and
$p(k)\le e^{3\sqrt k}$.
\end{lemma}
\begin{proof}
Orient a longest line as the first row.  Put $b=\lambda_2$ and
$h=\lambda'_1$.  For $k\ge1$, $b(h-1)\ge k$.
If $h-1\ge\sqrt k$, the diagram contains a hook with row and
column length $h$, because $L\ge h$.  That hook has
$\binom{2h-2}{h-1}\ge2^{h-1}$ tableaux.
Otherwise $b>\sqrt k$ and the diagram contains a two-row rectangle
of width $b$.  Its Catalan number is at least $2^{b-1}$: encode
the $2^{b-1}$ compositions of $b$ as distinct Dyck paths obtained
by concatenating pyramids of those lengths.  In this second case
$b\ge2$, so $b-1\ge b/2>\sqrt k/2$.
Subdiagram extension proves \eqref{l:degree-sqrt-deficit-equation};
the case $k=0$ is immediate.
Finally sum over deficits using at most $2p(k)$ shapes:
\[
 \sum_{\lambda\vdash n}D_\lambda^{-32}
 \le2+\sum_{k\ge1}2e^{3\sqrt k}e^{-16(\log2)\sqrt k}<\infty.
\]
\end{proof}

\subsection{Further direct fixed-power constructions}
\label{l:degree-fixed-powers}

Two final counts retain the larger powers that arise when only
especially rough estimates are used.  They are independent
specializations of the tableau constructions, not deductions from
an already proved smaller inverse power.

For power $30$ and $1\le k\le n/4$, choose all $k$ lower labels
from $\{2k+1,\ldots,n\}$.  Fill them in one fixed standard order
and put the other labels increasingly in the first row.  Since the
first $2k$ labels then lie in the top row, every lower predecessor
condition holds.  This gives
\begin{equation}\label{l:degree-thirty-prefill}
 D_\lambda\ge\binom{n-2k}k
       \ge(n/(2k))^k\ge2^k.
\end{equation}
The $2\cdot2^k$ shape count supplies the summable majorant
$2\cdot2^{-29k}$ and fixed-$k$ convergence.
For $k>n/4$, put $b=\lfloor\sqrt n/4\rfloor$ and retain the
first row $L\ge\sqrt n$ together with $b$ lower boxes.
Applying the same over-reserving construction to this subdiagram
gives $D_\lambda\ge\binom{L-b}b\ge3^b$.
Since $30\log3/4>3$, its inverse thirtieth power dominates
$p(n)\le e^{3\sqrt n}$.  Hence $Z_{30}(n)\to0$ and $C_{30}<\infty$.

For power $32$, a proof using the rough count $2^n$ separates
three ranges.  If $L\le n/8$, hooks give $(4/e)^n$, which is
enough because $32\log(4/e)>\log2$.
If $L>n/8$ and $k>n/3$, then $k>L/2$ and retaining
$\lfloor L/2\rfloor$ lower boxes gives
$D_\lambda\ge2^{\lfloor L/2\rfloor}$; its inverse thirty-second
power also dominates $2^n$.  For $1\le k\le n/3$, use
\[
 D_\lambda\ge\binom{n-k}k,\qquad
 2^{k+1}\left(\frac{k}{n-k}\right)^{32k}\le2^{1-31k},
\]
with fixed-$k$ convergence.  This proves the same complete
inverse-thirty-second assertion without using the square-root
partition bound.

For power $6$, the hook ratio gives
$D_\lambda\ge e^{-1}\binom nk$ for $1\le k\le n/3$.
The near-row and near-column contribution is bounded by
\begin{equation}\label{l:degree-six-hook-sum}
 2e^6\sum_{1\le k\le n/3}[2(k/n)^6]^k=o(1).
\end{equation}
For $k\le\sqrt n$ compare with the geometric series of base
$2n^{-3}$; for $k>\sqrt n$ use $2(k/n)^6\le2/3^6<1$.
In the remainder $L\le2n/3$.  If $L\ge n/(4e)$, retain
$b=\lfloor L/2\rfloor$ lower boxes and apply the hook ratio to
the resulting diagram, obtaining $e^{-1}\binom{L+b}b\ge e^{-1}3^b$.
If $L<n/(4e)$, hooks give $(n/(2eL))^n>2^n$.
These exponential bounds finish the sum with $p(n)=e^{O(\sqrt n)}$.

\begin{proof}[A direct fourth-power proof]
\phantomsection\label{l:degree-direct-four}
For power $4$, one may also use the near-row majorant
$2\cdot2^{-3k}$ on $k\le n/3$, the middle bound
$2^L/(L+1)$ on $n/8\le L<2n/3$, and the hook bound $(4/e)^n$
on $L<n/8$.  This is a different use of the fourth power from
the tail-adapted diagonal argument in Lemma~\ref{l:degree-diagonal-four}.
It uses the subexponential total partition count in its last two ranges.
\end{proof}

The sharper quantitative estimate
\[
 \sum_{\lambda\vdash n}D_\lambda^{-u}=2+O(n^{-u})
 \qquad(u>0\text{ fixed})
\]
is due to Liebeck and Shalev \cite[Theorem~2.6]{liebeck-shalev2004}.
The arguments above supply the boundedness and vanishing assertions
used here without that rate.

The corresponding alternating-group statement is
\begin{equation}\label{l:alternating-degree-sum}
 \sum_{\chi\in\Irr(A_n)}\chi(1)^{-u}=1+o(1)
 \qquad(u>0\text{ fixed}).
\end{equation}
Here $A_n$ is the subgroup of even permutations.  For completeness,
the index-two restriction rule makes the two vanishing assertions
equivalent.  If $V$ is irreducible for $S_n$, Frobenius reciprocity and
$\operatorname{Ind}_{A_n}^{S_n}\operatorname{Res}_{A_n}^{S_n}V
 \cong V\oplus(V\otimes\mathrm{sgn})$
show that the endomorphism algebra of its restriction has dimension
one or two.  Thus the restriction is irreducible or is a sum of two
distinct irreducibles.  In the latter case an odd permutation exchanges
the two summands, since otherwise either would be an $S_n$-invariant
proper subspace.  Both therefore have dimension $\dim V/2$.
Every irreducible of $A_n$ occurs in some restriction, by induction and
Frobenius reciprocity; it occurs in restrictions of at most two
$S_n$-types, since its induced representation has dimension twice its
own and each such type has at least its dimension.
For $n\ge5$, write $Z_u^A(n)$ for the sum over the nontrivial
$A_n$-types.  The only $S_n$-types restricting to the trivial type
are the trivial and sign representations, so the preceding counts give
\[
 Z_u(n)\le2Z_u^A(n),\qquad
 Z_u^A(n)\le2^{u+1}Z_u(n).
\]
This proves the equivalence and \eqref{l:alternating-degree-sum}, while
retaining the sharper cited symmetric-group rate separately.
